% Batch Normalization Tutorial
% Standalone LaTeX document

\documentclass[11pt,a4paper]{article}

% ------------------------------------------------------------
% Packages
% ------------------------------------------------------------
\usepackage[utf8]{inputenc}
\usepackage[T1]{fontenc}
\usepackage{lmodern}
\usepackage{amsmath,amssymb,mathtools}
\usepackage{geometry}
\usepackage{xcolor}
\usepackage{booktabs}
\usepackage{array}
\usepackage{enumitem}
\usepackage{tcolorbox}
\usepackage{tikz}
\usepackage{hyperref}
\usepackage{microtype}

% ------------------------------------------------------------
% Page setup
% ------------------------------------------------------------
\geometry{
	top=1in,
	bottom=1in,
	left=0.85in,
	right=0.85in
}

\hypersetup{
	colorlinks=true,
	linkcolor=blue!60!black,
	urlcolor=blue!60!black
}

\setlength{\parindent}{0pt}
\setlength{\parskip}{6pt}

% ------------------------------------------------------------
% Colors
% ------------------------------------------------------------
\definecolor{titleblue}{RGB}{20,75,125}
\definecolor{lightblue}{RGB}{235,245,255}
\definecolor{lightgreen}{RGB}{235,250,240}
\definecolor{lightorange}{RGB}{255,246,225}
\definecolor{lightpurple}{RGB}{245,240,255}
\definecolor{lightred}{RGB}{255,240,240}

% ------------------------------------------------------------
% Custom boxes
% ------------------------------------------------------------
\newtcolorbox{keyidea}{
	colback=lightblue,
	colframe=titleblue,
	boxrule=0.8pt,
	arc=3pt,
	title=\textbf{Key Idea}
}

\newtcolorbox{important}{
	colback=lightorange,
	colframe=orange!70!black,
	boxrule=0.8pt,
	arc=3pt,
	title=\textbf{Important}
}

\newtcolorbox{intuition}{
	colback=lightgreen,
	colframe=green!50!black,
	boxrule=0.8pt,
	arc=3pt,
	title=\textbf{Intuition}
}

\newtcolorbox{summarybox}{
	colback=lightpurple,
	colframe=purple!60!black,
	boxrule=0.8pt,
	arc=3pt,
	title=\textbf{Quick Summary}
}

% ------------------------------------------------------------
% Title
% ------------------------------------------------------------
\title{
	\vspace{-1cm}
	\textbf{\color{titleblue}Batch Normalization}\\[4pt]
	\large A Short Tutorial for Deep Learning Students
}
\author{}
\date{}

\begin{document}
	
	\maketitle
	
	\begin{keyidea}
		\textbf{Batch Normalization (BatchNorm)} is a technique used in deep
		neural networks to normalize intermediate activations within a mini-batch.
		It generally makes training faster, more stable, and less sensitive to
		the choice of initialization.
	\end{keyidea}
	
	\section{Why Batch Normalization?}
	
	Consider a neural network layer that receives an input $x$ and produces
	an intermediate value
	
	\begin{equation}
		z = Wx+b.
	\end{equation}
	
	Here, $W$ represents the weights and $b$ represents the bias.
	
	During training, the weights of previous layers continuously change.
	Consequently, the distribution and scale of the intermediate activations
	can also change. Batch Normalization helps control these activations by
	normalizing them using statistics computed from a mini-batch.
	
	\begin{intuition}
		Imagine that one layer produces values that are sometimes very small and
		sometimes very large. The next layer has to continuously adapt to these
		changes. BatchNorm keeps the activations in a more controlled form before
		passing them forward.
	\end{intuition}
	
	\section{Input to Batch Normalization}
	
	Suppose a mini-batch contains $m$ training examples and the layer produces
	$d$ features. The activations can be represented by the matrix
	
	\begin{equation}
		Z =
		\begin{bmatrix}
			z_{11} & z_{12} & \cdots & z_{1d}\\
			z_{21} & z_{22} & \cdots & z_{2d}\\
			\vdots & \vdots & \ddots & \vdots\\
			z_{m1} & z_{m2} & \cdots & z_{md}
		\end{bmatrix}.
	\end{equation}
	
	The rows correspond to training examples and the columns correspond to
	features.
	
	For each feature $j$, BatchNorm calculates its statistics across the
	$m$ examples in the mini-batch.
	
	\section{Step 1: Compute the Batch Mean}
	
	For feature $j$, the mean is calculated as
	
	\begin{equation}
		\mu_j =
		\frac{1}{m}
		\sum_{i=1}^{m} z_{ij}.
	\end{equation}
	
	Thus, every feature has its own mean.
	
	For example, if the layer has three features, we calculate
	
	\[
	\mu_1,\qquad \mu_2,\qquad \mu_3.
	\]
	
	\section{Step 2: Compute the Batch Variance}
	
	The variance for feature $j$ is
	
	\begin{equation}
		\sigma_j^2 =
		\frac{1}{m}
		\sum_{i=1}^{m}
		(z_{ij}-\mu_j)^2.
	\end{equation}
	
	The mean tells us where the activations are centered, while the variance
	tells us how much they are spread out.
	
	\begin{important}
		The mean and variance are computed separately for each feature, using the
		examples in the current mini-batch.
	\end{important}
	
	\section{Step 3: Normalize the Activations}
	
	Once the mean and variance are available, each activation is normalized as
	
	\begin{equation}
		\hat{z}_{ij}
		=
		\frac{z_{ij}-\mu_j}
		{\sqrt{\sigma_j^2+\epsilon}},
	\end{equation}
	
	where $\epsilon$ is a small positive number added for numerical stability.
	
	The purpose of this operation is to approximately obtain
	
	\begin{equation}
		\mathbb{E}[\hat{z}_j]\approx 0,
		\qquad
		\operatorname{Var}(\hat{z}_j)\approx 1.
	\end{equation}
	
	The normalized activations can be written as
	
	\begin{equation}
		\hat{Z} =
		\begin{bmatrix}
			\hat{z}_{11} & \hat{z}_{12} & \cdots & \hat{z}_{1d}\\
			\hat{z}_{21} & \hat{z}_{22} & \cdots & \hat{z}_{2d}\\
			\vdots & \vdots & \ddots & \vdots\\
			\hat{z}_{m1} & \hat{z}_{m2} & \cdots & \hat{z}_{md}
		\end{bmatrix}.
	\end{equation}
	
	\section{Step 4: Scale and Shift}
	
	If BatchNorm always forced every feature to have mean $0$ and variance
	$1$, it could unnecessarily restrict what the network can learn.
	
	Therefore, BatchNorm introduces two learnable parameters:
	
	\begin{itemize}[leftmargin=1.5em]
		\item $\gamma_j$: learnable scale parameter,
		\item $\beta_j$: learnable shift parameter.
	\end{itemize}
	
	The final BatchNorm output is
	
	\begin{equation}
		\boxed{
			y_{ij} = \gamma_j\hat{z}_{ij}+\beta_j
		}
	\end{equation}
	
	The parameters $\gamma_j$ and $\beta_j$ are learned during training using
	backpropagation.
	
	\begin{intuition}
		Normalization provides a convenient standard form, while $\gamma$ and
		$\beta$ give the network the freedom to learn the scale and position that
		are useful for the task.
	\end{intuition}
	
	\section{Complete Batch Normalization Formula}
	
	The complete process can be summarized as follows.
	
	\subsection*{1. Compute the mean}
	
	\begin{equation}
		\mu_j =
		\frac{1}{m}\sum_{i=1}^{m}z_{ij}.
	\end{equation}
	
	\subsection*{2. Compute the variance}
	
	\begin{equation}
		\sigma_j^2 =
		\frac{1}{m}
		\sum_{i=1}^{m}(z_{ij}-\mu_j)^2.
	\end{equation}
	
	\subsection*{3. Normalize}
	
	\begin{equation}
		\hat{z}_{ij}
		=
		\frac{z_{ij}-\mu_j}
		{\sqrt{\sigma_j^2+\epsilon}}.
	\end{equation}
	
	\subsection*{4. Scale and shift}
	
	\begin{equation}
		\boxed{
			y_{ij}=\gamma_j\hat{z}_{ij}+\beta_j
		}
	\end{equation}
	
	Thus,
	
	\begin{equation}
		\boxed{
			z
			\longrightarrow
			\text{Mean/Variance}
			\longrightarrow
			\text{Normalize}
			\longrightarrow
			\text{Scale and Shift}
			\longrightarrow
			y
		}
	\end{equation}
	
	\section{BatchNorm in a Neural Network}
	
	A common arrangement in a neural network is
	
	\begin{equation}
		\boxed{
			\text{Input}
			\rightarrow
			\text{Linear Layer}
			\rightarrow
			\text{BatchNorm}
			\rightarrow
			\text{ReLU}
			\rightarrow
			\text{Next Layer}
		}
	\end{equation}
	
	For example,
	
	\[
	x
	\rightarrow
	(Wx+b)
	\rightarrow
	\operatorname{BatchNorm}
	\rightarrow
	\operatorname{ReLU}.
	\]
	
	The BatchNorm operation is therefore performed on the intermediate
	activations before they are passed to the activation function and the
	subsequent layer.
	
	\section{Training Versus Inference}
	
	One of the most important concepts in BatchNorm is that the statistics
	used during training and inference are different.
	
	\subsection{During Training}
	
	During training, the mean and variance are calculated from the current
	mini-batch:
	
	\begin{equation}
		\mu_j=\mu_{j,\mathrm{batch}},
		\qquad
		\sigma_j^2=\sigma_{j,\mathrm{batch}}^2.
	\end{equation}
	
	At the same time, BatchNorm maintains running estimates of these
	statistics, commonly using an exponential moving average.
	
	\subsection{During Inference}
	
	During testing or inference, the model should not depend on the particular
	examples present in a mini-batch. Therefore, the stored running estimates
	are used:
	
	\begin{equation}
		y =
		\gamma
		\frac{z-\mu_{\mathrm{running}}}
		{\sqrt{\sigma_{\mathrm{running}}^2+\epsilon}}
		+\beta.
	\end{equation}
	
	The learned $\gamma$ and $\beta$ parameters are used as fixed parameters
	during inference.
	
	\begin{important}
		A common mistake is to use the statistics of the current test example
		during inference. BatchNorm normally uses the running mean and variance
		estimated during training.
	\end{important}
	
	\section{Benefits of Batch Normalization}
	
	BatchNorm can provide several practical benefits:
	
	\begin{enumerate}[leftmargin=1.8em]
		\item \textbf{Faster convergence:}
		Training often converges more quickly.
		
		\item \textbf{More stable training:}
		The scale of intermediate activations is controlled.
		
		\item \textbf{Higher learning rates:}
		The network can often tolerate relatively larger learning rates.
		
		\item \textbf{Regularization effect:}
		The randomness introduced by mini-batch statistics can sometimes
		reduce overfitting.
		
		\item \textbf{Reduced sensitivity to initialization:}
		Training can become less sensitive to the initial values of weights.
	\end{enumerate}
	
	\section{What Does BatchNorm Actually Normalize?}
	
	This is an important conceptual distinction.
	
	\begin{summarybox}
		\[
		\boxed{\text{BatchNorm normalizes activations, not weights.}}
		\]
		
		For a layer
		
		\[
		z=Wx+b,
		\]
		
		BatchNorm operates on $z$ (or the corresponding layer activations),
		rather than directly normalizing $W$.
	\end{summarybox}
	
	This is different from techniques such as weight decay or $L_2$
	regularization, which directly affect the magnitude of the weights.
	
	\section{Simple Numerical Example}
	
	Consider one feature whose activations in a mini-batch are
	
	\[
	z=[2,\;4,\;6].
	\]
	
	The batch mean is
	
	\begin{equation}
		\mu=\frac{2+4+6}{3}=4.
	\end{equation}
	
	The variance is
	
	\begin{equation}
		\sigma^2
		=
		\frac{(2-4)^2+(4-4)^2+(6-4)^2}{3}
		=
		\frac{8}{3}.
	\end{equation}
	
	Ignoring the very small $\epsilon$ for simplicity, the normalized values
	are
	
	\begin{equation}
		\hat{z}
		=
		\left[
		\frac{-2}{\sqrt{8/3}},
		\frac{0}{\sqrt{8/3}},
		\frac{2}{\sqrt{8/3}}
		\right].
	\end{equation}
	
	Therefore, the values are centered around zero and have approximately
	unit variance.
	
	If the learned parameters are, for example,
	
	\[
	\gamma=2,
	\qquad
	\beta=1,
	\]
	
	then
	
	\begin{equation}
		y=2\hat{z}+1.
	\end{equation}
	
	Thus, normalization is followed by a learned scale and shift.
	
	\section{BatchNorm: Four Steps to Remember}
	
	Students can remember BatchNorm using the following four-step sequence:
	
	\begin{enumerate}
		\item \textbf{Mean} -- calculate the batch mean.
		\item \textbf{Variance} -- calculate the batch variance.
		\item \textbf{Normalize} -- obtain approximately zero mean and unit variance.
		\item \textbf{Scale + Shift} -- use learnable $\gamma$ and $\beta$.
	\end{enumerate}
	
	In compact form,
	
	\begin{equation}
		\boxed{
			\text{Mean}
			\rightarrow
			\text{Variance}
			\rightarrow
			\text{Normalize}
			\rightarrow
			\text{Scale + Shift}
		}
	\end{equation}
	
	\section{Key Points for Examination}
	
	\begin{itemize}[leftmargin=1.5em]
		\item BatchNorm normalizes intermediate activations.
		\item Mean and variance are calculated for each feature across a
		mini-batch during training.
		\item $\epsilon$ is added for numerical stability.
		\item $\gamma$ and $\beta$ are learnable parameters.
		\item $\gamma$ controls scale and $\beta$ controls shift.
		\item During inference, running estimates of mean and variance are used.
		\item BatchNorm can improve training stability and convergence.
		\item BatchNorm is not the same as weight regularization.
	\end{itemize}
	
	\section{One-Line Definition}
	
	\begin{tcolorbox}[
		colback=lightblue,
		colframe=titleblue,
		boxrule=1pt,
		arc=4pt
		]
		\centering
		\large
		\textbf{
			Batch Normalization normalizes the intermediate activations of a
			neural network using mini-batch statistics and then restores useful
			scale and shift through learnable parameters $\gamma$ and $\beta$.
		}
	\end{tcolorbox}
	
\end{document}
